\documentclass[11pt, letterpaper]{article}

\usepackage[margin=1in, headheight=22pt, headsep=12pt]{geometry}
\usepackage{titlesec}
\usepackage{enumitem}
\usepackage{xcolor}
\usepackage{hyperref}
\usepackage{fancyhdr}
\usepackage{microtype}
\usepackage{parskip}
\usepackage{booktabs}
\usepackage{sectsty}
\usepackage{amssymb}

% --- Color Definitions ---
\definecolor{accent}{HTML}{2C3E6B}
\definecolor{subaccent}{HTML}{5B7BA5}
\definecolor{rulecolor}{HTML}{CCCCCC}
\definecolor{darktext}{HTML}{1A1A1A}
\definecolor{greenzone}{HTML}{2E7D32}
\definecolor{yellowzone}{HTML}{F9A825}
\definecolor{redzone}{HTML}{C62828}
\definecolor{grayzone}{HTML}{616161}
\definecolor{lightbg}{HTML}{F5F5F5}

% --- Hyperlink Setup ---
\hypersetup{
    colorlinks=true,
    linkcolor=accent,
    urlcolor=subaccent,
    pdftitle={Week 1: Early Jazz - Quiz Study Guide},
    pdfauthor={MUS 262},
}

% --- Section Formatting ---
\titleformat{\section}
    {\Large\bfseries\color{accent}}
    {\thesection}{1em}{}
    [\vspace{-0.5em}\textcolor{rulecolor}{\rule{\textwidth}{0.4pt}}]

\titleformat{\subsection}
    {\large\bfseries\color{subaccent}}
    {\thesubsection}{1em}{}

\titleformat{\subsubsection}
    {\normalsize\bfseries\color{accent}}
    {\thesubsubsection}{1em}{}

\titlespacing*{\section}{0pt}{1.5em}{0.75em}
\titlespacing*{\subsection}{0pt}{1em}{0.5em}
\titlespacing*{\subsubsection}{0pt}{0.75em}{0.4em}

% --- Header/Footer ---
\pagestyle{fancy}
\fancyhf{}
\renewcommand{\headrulewidth}{0.4pt}
\renewcommand{\headrule}{\hbox to\headwidth{\color{rulecolor}\leaders\hrule height \headrulewidth\hfill}}
\fancyhead[L]{\small\color{gray}MUS 262 --- Week 1: Early Jazz}
\fancyhead[R]{\small\color{gray}Quiz Study Guide}
\fancyfoot[C]{\small\color{gray}\thepage}

% --- List Formatting ---
\setlist[itemize]{leftmargin=1.5em, itemsep=3pt, parsep=0pt, topsep=4pt}
\setlist[itemize,2]{leftmargin=1.5em, itemsep=2pt}
\setlist[enumerate]{leftmargin=1.5em, itemsep=3pt, parsep=0pt, topsep=4pt}

% --- Custom Commands ---
\newcommand{\tracker}[2]{$\square$ \textbf{#1} #2}
\newcommand{\listenbox}[1]{\vspace{0.5em}\noindent\fcolorbox{rulecolor}{lightbg}{\parbox{\dimexpr\textwidth-2\fboxsep-2\fboxrule}{\small\textbf{Listen For:} #1}}\vspace{0.5em}}

% --- Body Text ---
\color{darktext}

\begin{document}

% --- Title Block ---
\thispagestyle{empty}
\begin{center}
    \vspace*{0.5cm}
    {\Huge\bfseries\color{accent} Week 1: Early Jazz}\\[0.5em]
    {\Large\color{subaccent} Quiz Study Guide}\\[1em]
    \textcolor{rulecolor}{\rule{0.6\textwidth}{0.8pt}}\\[0.8em]
    {\color{gray}\small MUS/AAS 262 --- Jazz History: Many Sounds, Many Voices\\[0.3em]
    New Orleans Jazz (1900--1928)}
    \vspace{1cm}
\end{center}

\tableofcontents
\newpage

% ============================================================
\section{Knowledge Tracker}
% ============================================================

\textit{Move items between sections as you study. The goal is to get everything into Green.}

\subsection*{\textcolor{greenzone}{Green --- Know From Memory}}
\begin{itemize}[label=$\square$]
    \item Louis Armstrong = ``First Great Soloist'' (Schuller)
    \item Jazz emphasizes beats 2 and 4 (classical = 1 and 3)
\end{itemize}

\subsection*{\textcolor{yellowzone}{Yellow --- Partial / Need Notes}}
\begin{itemize}[label=$\square$]
    \item Freddie Keppard refused first recording offer (didn't want people to copy him)
    \item Blues scale = pentatonic + tritone; blue notes = flat 3, flat 5, flat 7
    \item Heterophony = no one plays the same melodic line at the same time
    \item ODJB = Original Dixieland Jazz Band, first jazz recording, 1917, all-white
\end{itemize}

\subsection*{\textcolor{redzone}{Red --- Didn't Know / Got Wrong}}
\begin{itemize}[label=$\square$]
    \item Congo Square = where \textit{enslaved people} in 19th-century New Orleans congregated to make music (preserved African traditions, predates jazz)
    \item Bessie Smith = ``Empress of the Blues''; bending and sliding into notes influenced instrumentalists
    \item Jelly Roll Morton = Creole (mixed Black + European); claimed to have invented jazz; Red Hot Peppers
    \item 12-bar blues = A, A, B (4+4+4 bars) --- NOT AABA (that is 32-bar form, different thing)
    \item Creole = mixed-race (Black + European), classically trained, barred from white orchestras
\end{itemize}

\subsection*{\textcolor{grayzone}{Not Yet Tested}}
\begin{itemize}[label=$\square$]
    \item Scott Joplin = ``King of Ragtime''; Maple Leaf Rag (1899 sheet music, 1916 piano roll)
    \item W.C. Handy = ``Father of the Blues''; St. Louis Blues has habanera/tango rhythm (Spanish tinge)
    \item Armstrong invented scat singing (1926, ``Heebie Jeebies'')
    \item Armstrong: vibrato only at END of notes (revolutionary technique)
    \item ``West End Blues'' (1928) = unaccompanied trumpet cadenza opening; written by King Oliver
    \item Hot Five / Hot Seven = Armstrong's landmark recording groups (1925--1928)
    \item Morton's ``Dead Man Blues'' (1926) = spoken dialogue intro, funeral march, Red Hot Peppers
    \item Bessie Smith ``Backwater Blues'' (1927) = dual call-and-response, James P. Johnson on piano
    \item ODJB ``Dixie Jass Band One-Step'' = marching band influence, rehearsed, minimal improv
    \item Handy = formally trained, first to transcribe blues into written notation
    \item Call and response = fundamental jazz element; ``trading eights'' in instrumental jazz
    \item Swing rhythm = dotted eighth + sixteenth, or triplet with tie on first two notes
    \item Baraka quote: ``Blues could not exist if the African captives had not become American captives''
    \item New Orleans instrumentation = trumpet/cornet, clarinet, trombone (reduction of marching band)
    \item Ragtime = syncopated melodies, steady bass, NO improvisation, precursor to jazz
    \item 32-bar AABA form = popular song form (B section = bridge) --- different from 12-bar blues
\end{itemize}

\newpage

% ============================================================
\section{The Era: Early Jazz / New Orleans Jazz (1900--1928)}
% ============================================================

\subsection{What This Movement Contributes}
\begin{enumerate}
    \item \textbf{Birth of Jazz as a Genre} --- emerged in New Orleans around 1900 as fusion of African rhythms, blues, ragtime, work songs, European harmony
    \item \textbf{Collective Improvisation} --- multiple instruments improvising simultaneously (heterophony)
    \item \textbf{Louis Armstrong's Revolution} --- transformed jazz from ensemble music to soloist's art form; invented scat singing
    \item \textbf{Blues Form Codification} --- 12-bar blues structure (A, A, B) became foundational
    \item \textbf{Commercial Viability} --- first jazz recordings spread the music beyond New Orleans
\end{enumerate}

\subsection{Era Characteristics}
\begin{itemize}
    \item \textbf{Rhythmic Feel:} Swing rhythm (dotted eighth + sixteenth, or triplet with tie); beats 2 and 4
    \item \textbf{Ensemble Style:} Small groups; trumpet/cornet, clarinet, trombone as core (reduction of marching band)
    \item \textbf{Improvisation:} Collective improvisation; heterophony
    \item \textbf{Blues Influence:} Blues scale (pentatonic + tritone), blue notes (flat 3, flat 5, flat 7), call and response
    \item \textbf{Recording Tech:} Limited to $\sim$3 minutes per disc side
    \item \textbf{Sound:} Instrumental playing mimics vocal techniques (bends, moans, vibrato)
\end{itemize}

\subsection{Key Concepts}
\begin{itemize}
    \item \textbf{Congo Square:} Where enslaved people in 19th-century New Orleans congregated to make music
    \item \textbf{Heterophony:} No one plays the same melodic line at the same time
    \item \textbf{Blues Structure:} A, A, B form; each section = 4 bars (12 bars total)
    \item \textbf{Blues Scale:} Pentatonic scale + tritone (the ``Devil's Interval'')
    \item \textbf{Call and Response:} Biggest element of jazz; became ``trading eights'' in instrumental jazz
    \item \textbf{Swing Rhythm:} Dotted eighth + sixteenth, or triplet with tie on first two notes
    \item \textbf{Creole:} Mixed-race (Black + European), classically trained, barred from white orchestras
    \item \textbf{Baraka quote:} ``Blues could not exist if the African captives had not become American captives''
\end{itemize}

\newpage

% ============================================================
\section{Artists and Songs}
% ============================================================

% --- JOPLIN ---
\subsection{Scott Joplin --- ``Maple Leaf Rag'' (1916, Piano Roll)}

\textbf{Era:} Ragtime (precursor to jazz), late 1890s--1910s

\subsubsection*{Key Facts}
\begin{enumerate}
    \item ``King of Ragtime''
    \item Published as sheet music in 1899 --- one of first to sell over 1 million copies
    \item 1916 recording is a piano roll, not live --- played back mechanically
    \item Ragtime = syncopated melodies against steady bass; direct precursor to jazz
    \item NO improvisation --- fully composed and notated
    \item Born 1868 (right after Emancipation), died 1917
\end{enumerate}

\textbf{Instrumentation:} Solo piano (mechanical piano roll)\\
\textbf{Song Form:} Multi-strain (AABBACCDD); NOT 12-bar blues or 32-bar AABA

\listenbox{Syncopated melody over steady ``oom-pah'' bass, stiff/mechanical rhythm, no swing feel, bright/playful, clear repeated sections, no improvisation}

\bigskip
\textcolor{rulecolor}{\rule{\textwidth}{0.2pt}}

% --- HANDY ---
\subsection{W.C. Handy --- ``St. Louis Blues'' (1922)}

\textbf{Era:} Early Jazz / Blues (1910s--1920s)

\subsubsection*{Key Facts}
\begin{enumerate}
    \item ``Father of the Blues''
    \item Composed 1914, this recording 1922
    \item Formally trained; first to transcribe and publish blues in written notation
    \item Habanera (tango) rhythm in middle section --- Cuban influence (``Spanish tinge'')
    \item Bandleader, not rural folk musician --- professional, literate Black musician class
    \item Pianist and trumpet player
\end{enumerate}

\textbf{Instrumentation:} Small ensemble --- trumpet (muted and open), clarinet, piano, bass, drums\\
\textbf{Song Form:} Blues-based but complex; sections switch between swung/straight, minor/major; tango section in middle

\listenbox{Highly composed, very little improvisation, polyphony to monophony, swung vs.\ straight switches, minor/major shifts, muted trumpet timbre, tango rhythm}

\bigskip
\textcolor{rulecolor}{\rule{\textwidth}{0.2pt}}

% --- ODJB ---
\subsection{Original Dixieland Jazz Band --- ``Dixie Jass Band One-Step'' (1917)}

\textbf{Era:} Early Jazz / New Orleans Jazz (1917)

\subsubsection*{Key Facts}
\begin{enumerate}
    \item First commercially released jazz recording (1917, Victor Records)
    \item All-white band --- brought jazz to mainstream white audiences
    \item Freddie Keppard (Black cornet player) refused first chance to record --- feared copying
    \item Clearly rehearsed with minimal improvisation --- strong marching band influence
    \item From New Orleans, breakthrough in Chicago and New York
    \item Originally spelled ``J-A-S-S'' (later ``J-A-Z-Z'')
\end{enumerate}

\textbf{Instrumentation:} Trumpet/cornet, clarinet, trombone, piano, drums (marching band tradition)\\
\textbf{Song Form:} Multi-strain ragtime-influenced; highly rehearsed

\listenbox{Marching band feel, stiff/military, clearly rehearsed, minimal improv, trumpet/clarinet/trombone lines, sounds like ragtime, no real soloists}

\bigskip
\textcolor{rulecolor}{\rule{\textwidth}{0.2pt}}

% --- BESSIE SMITH ---
\subsection{Bessie Smith --- ``Backwater Blues'' (1927)}

\textbf{Era:} Classic Blues / Early Jazz (1920s)

\subsubsection*{Key Facts}
\begin{enumerate}
    \item ``Empress of the Blues'' (1894--1937) --- most commercially successful blues artist of 1920s
    \item Inspired by Great Mississippi Flood of 1927
    \item Two simultaneous call-and-responses --- vocal/piano AND within vocal phrasing
    \item Bending and sliding into notes --- pattern instrumentalists later imitated
    \item Columbia Records; James P. Johnson on piano
    \item Darker timbre, extensive vibrato
\end{enumerate}

\textbf{Instrumentation:} Voice (lead) + Piano (James P. Johnson)\\
\textbf{Song Form:} 12-bar blues (A, A, B); each section = 4 bars

\listenbox{Dual call-and-response, vocal bending/sliding (blue notes), dark timbre, heavy vibrato, piano fills between vocal phrases, emotional storytelling}

\bigskip
\textcolor{rulecolor}{\rule{\textwidth}{0.2pt}}

% --- MORTON ---
\subsection{Jelly Roll Morton --- ``Dead Man Blues'' (1926)}

\textbf{Era:} Early Jazz / New Orleans Jazz (1920s)

\subsubsection*{Key Facts}
\begin{enumerate}
    \item Born Ferdinand Joseph LaMothe (1890--1941); Creole --- mixed-race, classically trained, barred from white orchestras
    \item Claimed to have invented jazz
    \item Opens with spoken dialogue mimicking African American speech
    \item Recorded with Red Hot Peppers for Victor Records
    \item Highly arranged yet swinging
    \item Nickname from playing piano at houses of prostitution
    \item Pioneering jazz composer/arranger
\end{enumerate}

\textbf{Instrumentation:} New Orleans ensemble --- trumpet, clarinet, trombone, piano, banjo/guitar, drums, tuba/bass\\
\textbf{Song Form:} New Orleans funeral march; arranged sections with space for improvisation

\listenbox{Spoken intro dialogue, funeral march / second line rhythm, collective improvisation, ragtime influence / stiff rhythm, clarinet and trumpet solos, more composed than improvised}

\bigskip
\textcolor{rulecolor}{\rule{\textwidth}{0.2pt}}

% --- ARMSTRONG ---
\subsection{Louis Armstrong --- ``West End Blues'' (1928)}

\textbf{Era:} Early Jazz / New Orleans Jazz (late 1920s)

\subsubsection*{Key Facts}
\begin{enumerate}
    \item Opens with unaccompanied trumpet cadenza --- landmark in jazz history
    \item ``First Great Soloist'' (Schuller Ch. 3) --- transformed jazz from ensemble to soloist's art
    \item First to play vibrato only at END of note --- influenced every instrument since
    \item Hot Five and Hot Seven recordings (1925--1928) = most important in jazz history
    \item Written by King Oliver, Armstrong's mentor who brought him to Chicago
    \item Born 1901 New Orleans, died 1971; self-taught orphan
    \item Invented scat singing in 1926 with ``Heebie Jeebies''
\end{enumerate}

\textbf{Instrumentation:} Trumpet (Armstrong), clarinet, trombone, piano, banjo, drums (Hot Five combo)\\
\textbf{Song Form:} 12-bar blues; opens with unaccompanied trumpet cadenza

\listenbox{Iconic trumpet cadenza opening, strident/bright trumpet sound, vibrato at END of notes, improvised solo, trumpet mirrors vocal phrasing, blues structure}

\newpage

% ============================================================
\section{Listening Response Template}
% ============================================================

Name, composer, and year are PROVIDED on the quiz. Focus on what you HEAR. Aim for 5 observations:

\begin{enumerate}
    \item \textbf{Instrumentation/Ensemble} --- What instruments? Solo, combo, or large group? Who is soloist?
    \item \textbf{Era/Style} --- Ragtime, early blues, New Orleans jazz? What clues? (collective improv = early; virtuoso solo = Armstrong; stiff rhythm = ragtime)
    \item \textbf{Rhythm/Feel} --- Swinging or straight? Fast/slow? Beats 2 and 4? Syncopated? Second line?
    \item \textbf{Melody/Improvisation} --- Melody stated or improvised? Collective or solo improv? Bending, sliding, vibrato?
    \item \textbf{Texture/Form/Special} --- Dense or sparse? 12-bar blues, AABA, ragtime strains? Call and response? Special features (spoken intro, cadenza, muted trumpet, scat)?
    \item \textbf{Emotion/Context} --- What feeling? How does it connect to readings?
\end{enumerate}

\noindent\fcolorbox{rulecolor}{lightbg}{\parbox{\dimexpr\textwidth-2\fboxsep-2\fboxrule}{\small\textbf{Template:} ``This piece features [instruments], with [soloist] as the primary voice. The style is [era/style] based on [clues]. The form is [form], evidenced by [what you hear]. A notable feature is [detail], which connects to [course concept].''}}

% ============================================================
\section{General Understanding}
% ============================================================

\subsection{Movement's Place in History}
\begin{itemize}
    \item New Orleans birthplace: unique confluence of African, European, Caribbean, Creole cultures
    \item Congo Square: preserved African musical traditions that fed into jazz
    \item Ragtime (composed, syncopated) + Blues (vocal, 12-bar) = Early Jazz (improvised, ensemble)
    \item Jazz spread from New Orleans to Chicago and New York during Great Migration
    \item First recordings (1917): white musicians profited from Black innovation
\end{itemize}

\subsection{Cultural Context}
\begin{itemize}
    \item Minstrelsy: Armstrong criticized as ``old Negro'' for entertaining persona
    \item Creole class: classically trained but barred from white orchestras
    \item Harlem Renaissance: ``New Negro'' movement for Black artistic legitimacy
    \item Colorism: paper bag test at Cotton Club
\end{itemize}

\subsection{Musical Revolution}
\begin{itemize}
    \item African (polyrhythm, syncopation, call-response) + European (harmony, notation, phrase structure)
    \item Collective improvisation (heterophony) was precursor to solo improvisation
    \item Armstrong shifted jazz from ensemble to soloist's art
    \item Instrumental jazz mimicked vocal techniques from blues singers like Bessie Smith
\end{itemize}

\subsection{Key Tensions}
\begin{itemize}
    \item \textbf{Composed vs. Improvised:} ragtime composed; Armstrong introduced spontaneous improvisation
    \item \textbf{Entertainment vs. Art:} early jazz was dance music; tension with artistic legitimacy
    \item \textbf{Black Innovation, White Profit:} ODJB recorded first after Keppard refused
    \item \textbf{Preservation vs. Evolution:} Morton stayed 1920s style; Armstrong pushed forward
\end{itemize}

\end{document}
